% LSPV-3 REPORT::WRAPORBIT1::2026-10-12T12:37:13Z
% !TEX program = xelatex
\documentclass[11pt]{article}
\usepackage[margin=0.86in]{geometry}
\usepackage{fontspec,amsmath,amssymb,amsthm,booktabs,array,enumitem,xurl,hyperref,microtype}
\hypersetup{colorlinks=true,linkcolor=blue,urlcolor=blue,citecolor=blue}
\setlength{\emergencystretch}{3em}
\newtheorem{theorem}{Theorem}
\newtheorem{corollary}[theorem]{Corollary}
\newcommand{\SAT}{\mathsf{SAT}}
\newcommand{\EXACT}{\textbf{PROVED/EXACT}}
\newcommand{\KNOWN}{\textbf{KNOWN THEOREM}}
\newcommand{\COMP}{\textbf{COMPUTATIONAL}}
\newcommand{\CAND}{\textbf{CANDIDATE}}
\newcommand{\OPEN}{\textbf{OPEN}}
\newcommand{\RETIRED}{\textbf{FALSIFIED/RETIRED}}
\newcommand{\claim}[1]{\texttt{\detokenize{#1}}}

\title{LSPV-3 WRAPORBIT1:\\
Generated Wrapper Orbits Do Not Certify Universal Invariance}
\author{UNIVERSAL\_CORE\_REDTEAM}
\date{12 October 2026, 12:37 UTC}

\begin{document}
\sloppy
\maketitle

\begin{abstract}
P versus NP remains \OPEN.  Current audits correctly label finite wrapper
tests as computational, but the surviving full-domain semantic anti-sharing
target requires invariance under every equivalent realization.  This run
proves an exact non-certification theorem: even invariance under every finite
composition of a tested wrapper generator set establishes only invariance on
that generated orbit.  If one equivalent implementation lies outside the
orbit, an orbit indicator passes every generated test and fails universal
semantic invariance.  A deterministic audit over 77,750 generated programs
and 200 fresh equivalent implementations made 701,750 assertions with zero
failures.  The theorem retires only the inference from wrapper controls to
universal invariance; it does not refute a separately proved invariant.
\end{abstract}

\section{Strict result ledger}
\begin{center}\scriptsize
\begin{tabular}{>{\raggedright\arraybackslash}p{0.46\linewidth}>{\raggedright\arraybackslash}p{0.21\linewidth}>{\raggedright\arraybackslash}p{0.24\linewidth}}
\toprule
Claim & Status & Scope\\ \midrule
\path{TESTED-WRAPPER-MONOID-NONCERTIFICATION1} & \EXACT & Any declared program set and wrapper monoid\\
Concrete AST audit & \COMP & 701,750 assertions; zero failures\\
\path{FINITE-WRAPPER-CONTROLS-IMPLY-UNIVERSAL-INVARIANCE1} & \RETIRED & Logical inference without a coverage theorem\\
\path{FULL-DOMAIN-SEMANTIC-ANTI-SHARING1} & \OPEN & All exact uniform SAT deciders\\
P versus NP & \OPEN & No separation proved\\
\bottomrule
\end{tabular}
\end{center}

\section{Claim card and quantifiers attacked}
Let $\mathcal P$ be a set of deterministic program descriptions and
$\equiv$ extensional equality on a fixed input encoding.  Let
$G=\{g_1,\ldots,g_m\}$ be tested total semantics-preserving transformations,
for example COPY, DEAD, DELAY, INTERP, RECOMP, and DISPATCH.  Write
\[
 \langle G\rangle p
 =\{g_{i_k}\cdots g_{i_1}(p):k\ge0\}
\]
for the generated orbit.  The attacked inference is
\[
 \bigl[J(gp)=J(p)\text{ for all tested generators/compositions}\bigr]
 \Longrightarrow
 \bigl[J(q)=J(p)\text{ for every }q\equiv p\bigr].              \tag{1}
\]
The left side can be tested or proved within the declared wrapper grammar;
the right side quantifies over the complete equivalence class, including
independent circuits, branching programs, tables, source programs, and
interpreters not generated by that grammar.

\section{Exact orbit theorem}
\begin{theorem}[\claim{TESTED-WRAPPER-MONOID-NONCERTIFICATION1}]
Fix $p\in\mathcal P$ and a semantics-preserving transformation set $G$.  If
there exists $q\equiv p$ with $q\notin\langle G\rangle p$, then there exists a
decidable description functional $J:\mathcal P\to\{0,1\}$ satisfying
\[
 J(r)=1\quad\text{for every }r\in\langle G\rangle p,
 \qquad J(q)=0.
\]
Consequently $J$ is invariant under every tested generator wherever the
generator remains in the orbit, and under every finite composition of those
generators, but is not invariant on the semantic-equivalence class of $p$.
Thus (1) is false without the coverage premise
$[p]_{\equiv}=\langle G\rangle p$ or an independent invariance theorem.
\end{theorem}

\begin{proof}
Take $J$ to be the characteristic function of the generated orbit when orbit
membership is decidable in the declared wrapper grammar.  In the concrete
AST realization below, membership is the test that the protected base token
equals that of $p$ and every outer node has a generator label in $G$.
Every generated program passes, while $q$ fails by hypothesis.  More
generally, even without decidable full orbit membership, the finite set of
observations made by any bounded experiment can be separated from one fresh
equivalent description by a finite lookup predicate.  Hence observations on
the orbit do not logically entail the universal conclusion.
\end{proof}

\paragraph{Why the counterexample is legitimate but limited.}
The orbit indicator is intentionally adversarial.  It is not proposed as a
natural compatibility invariant.  Its role is to witness the missing
quantifier in a claimed certification rule.  If a candidate $J$ has a
separate mathematical proof of semantic invariance, this theorem does not
attack that proof.

\section{Concrete countermodel and reproducible audit}
The declared AST model stores a truth table, a base-implementation token, and
an outer wrapper chain.  The six standard generators preserve the token and
the table, although DELAY adds 17 charged steps.  Four fresh implementations
LUT, BDD, CIRCUIT, and ARITH use the same table with a different base token.
They are extensionally equivalent but outside the generated orbit.

The script
\path{experiments/review/wrapper_orbit_noncertification_audit.py} emitted
\path{results/raw/review/2026-10-12_wrapper_orbit_noncertification.json} in
0.535 seconds.  It checked:
\begin{itemize}[leftmargin=1.5em]
\item all 1,555 wrapper chains of depth at most four over the six generators;
\item ten families (2SAT, Horn, affine/Tseitin, bounded treewidth, parity,
independent products, padding, redundant clauses, signed clean controls, and
mixed SAT/UNSAT) under five encodings;
\item 77,750 generated programs and 200 fresh equivalent implementations;
\item semantic equality on eight inputs, orbit membership, wrapper depth, and
charged DELAY work; and
\item 701,750 assertions with zero failures.
\end{itemize}
Script SHA-256:
\texttt{845b94bae5e996a6d11f1819625e80542db6d0f1bdca0d0970de1a3f98d49b0c}.
Result SHA-256:
\texttt{606d3c8ae16d006bffcb8336077368ced3850b135991aef08df4dc46243a6af1}.
The audit is \COMP.  It validates only the concrete countermodel; the theorem
is symbolic and no asymptotic lower bound is inferred.

\section{Smallest repair}
Wrapper robustness must be established by one of the following, not by a
finite control suite:
\begin{enumerate}[leftmargin=1.6em]
\item prove that the tested transformation monoid is transitive on the exact
semantic-equivalence class (an extremely strong coverage theorem);
\item define $J$ directly on the extensional quotient and prove it is
well-defined, accepting that all exact SAT deciders have the same one-bit
function and implementation-specific information then disappears; or
\item give a covariant transport theorem with explicit source/target maps and
charge the compiler, validation, and reconstruction in one
\textsc{same\_realization} execution.
\end{enumerate}
Translation validation provides a useful software-verification analogy:
validating each produced translation can prove that particular transformation
correct, but it does not prove that a small set of transformations generates
all equivalent programs.  This is a \textbf{USEFUL STRUCTURAL ANALOGY
REQUIRING A NEW TRANSFER THEOREM}, not a SAT lower bound.

\section{Barrier, exponent, and endpoint audit}
The theorem survives trivial one-bit output and arbitrary polynomial
exponents because it concerns descriptions, not runtime.  DELAY wrappers
show that cost-sensitive functionals are not semantic; fresh LUT/circuit
implementations show that tested syntactic closure is not the full semantic
class.  Nonuniform per-length circuits enlarge the equivalence class further.

Rice's theorem is a \KNOWN{} guard: nontrivial extensional properties of
arbitrary partial computable functions are undecidable.  It does not prove
the orbit theorem and does not forbid a noncomputable mathematical invariant,
a restricted verified IR, or a proof-carrying transformation.  Pnueli,
Siegel, and Singerman's translation validation is directly transferable only
to an individually validated transformation; universal coverage is missing.
FLMV is directly transferable only after the same ordinary machine is proved
to satisfy its simultaneous time--space hypotheses.  No orbit argument
supplies that landing.

The result relativizes and algebrizes and constructs no large useful circuit
property, so it gives no Natural-Proofs escape.  It is a negative
certification theorem, not a nonrelativizing hinge.  A fixed-level proof,
communication, branching-program, VLSI, or control model supplies only a
model-specific orbit until an all-decider compiler and coverage theorem are
proved.

\section{Post-quantum technique ledger}
Schur squares, puncturing/shortening, quadratic-relation ranks, Hilbert
functions, syzygies, and graded Betti numbers are \textbf{FAILED TRANSFERS}
for universal wrapper coverage: they inspect a chosen representation and do
not show that every equivalent program exposes it.  Collision entropy gives
the exact inequality
$|\operatorname{im}L|\ge |A|^2/\operatorname{Coll}(L)$ for a fixed declared
map, an \textbf{EXACT OBJECT/THEOREM TRANSFER}; it cannot certify that the map
exists across the full equivalence class.  LWE-style external extraction and
special-soundness collision extraction are \textbf{ANALOGY/CANDIDATE}
requiring new source-matched compilers.  Random-oracle and cryptographic
binding assumptions are inadmissible.  Hidden algebraic structure can make a
fresh implementation easier, so it is not automatically a hardness source.

\section{CORPUS REVIEWED}
The complete current GitHub tree was surveyed.  This run read root
\path{main.tex}, \path{AGENTS.md}, \path{CODEX.md}, the complete
\path{.codex/skills/pnp-research/SKILL.md}, all current specifications under
\path{tasks/}, actual \path{main4.tex} and \path{main5.tex}, newest verified
MAIN SYNTHESIS-0023, and the newest LSPV-1 through LSPV-4 reports.  Materially
relevant history included CIRCEQUIV1, POINTINF1, REPDIRSUM1, ANTISHARE1,
MODELBRIDGE-0022/0024, ONETRANS1, APIFLAT1, TIGHTDELAY1, RT-0015,
COLLAPSEGATE1, DELAYOR1, and the recovered wrapper-variants audit.  The HEC,
VSC-HEC, G-PPC-HEC, CRS-PPC, and CCSPC causal branch, current governance,
and cited experiment/result artifacts were inspected.  Repository searches
covered finite wrappers, invariance, test suites, fresh wrappers, Rice-style
quotients, black-box extraction, and program equivalence.

Primary literature was checked before substantive work: Rice's original
semantic-property theorem, translation validation, relativization,
algebrization, Natural Proofs, and FLMV's SAT time--space lower bound.

\section{Predecessor write recovery}
The newest verified canonical LSPV-2 artifact is REPDIRSUM1 at 11:24 UTC.
No artifact attributable to the next scheduled LSPV-2 opportunity existed.
An exact later body was available with SHA-256 recorded above the current run,
so a nonblocking body-preserving recovery request was written and read back at
\path{research/review/outbox/redteam/2026-10-12_1237_LSPV2_RECOVERY_REQUEST.md}.
The evidence indicates a missing canonical write, not a repository-wide
permission failure; the LSPV-2 body was not reconstructed or published by
LSPV-3.

\section{DUPLICATION / CONTRADICTION AUDIT}
APIFLAT1 proves that named helper APIs can be erased by dispatcher flattening.
TIGHTDELAY1 and JOINTGAUGE1 show that time/space gauges change under equivalent
wrappers.  MODELBRIDGE-0022 proves black-box transcript extensionality;
POINTINF1 and RT-0015 block complete semantic canonicalization shortcuts.
Those results concern particular resource or quotient proposals.  WRAPORBIT1
adds the missing logical guard for wrapper \emph{testing}: even closure under
all compositions of the tested generators is only orbit-local.  It neither
duplicates nor contradicts a separately proved quotient or covariance
theorem.  Root MAIN already labels finite audits \COMP; this report explains
why that label is load-bearing.

\section{WHAT CHANGED RELATIVE TO ROOT MAIN}
Root MAIN's full-domain semantic anti-sharing target is unchanged and remains
\OPEN.  The LSPV-3 region now records that DIRECT/COPY/DEAD/DELAY/INTERP/
RECOMP (or any other finite generator suite) cannot certify its universal
wrapper-invariance premise.  Future positive claims must provide a coverage,
quotient, or covariance theorem before using finite controls as evidence for
the all-realization quantifier.

\section{ADOPTED / REJECTED / RETIRED PRIOR COMPONENTS}
\textbf{Adopted:} APIFLAT1, TIGHTDELAY1, JOINTGAUGE1, MODELBRIDGE-0022,
ANTISHARE1, exact wrapper charging, and strict \textsc{same\_realization}
discipline.  \textbf{Rejected:} extrapolation from a finite wrapper battery,
or from its generated monoid without a transitivity theorem, to every
semantically equivalent program.  \textbf{FALSIFIED/RETIRED:}
\claim{FINITE-WRAPPER-CONTROLS-IMPLY-UNIVERSAL-INVARIANCE1}.
\textbf{OPEN:} full-domain semantic anti-sharing, a source-matched compiler,
solver-independent simultaneous exponents, and a model-matched endpoint.

\section{STALE OR OVERSTATED CLAIMS FOUND}
Any statement that a candidate ``is wrapper invariant'' solely because it
passes the standard six controls is overstated.  The defensible label is
``passed these controls'' (\COMP) until a coverage/quotient/covariance theorem
is supplied.  Conversely, the existence of an adversarial orbit indicator
does not show that a separately proved mathematical invariant fails; such a
promotion would also be overstated.

\section{Precise handoff}
For \claim{FULL-DOMAIN-SEMANTIC-ANTI-SHARING1}, specify the exact program
equivalence relation and prove one of: (i) a transitive wrapper generator
coverage theorem; (ii) a nonconstant, well-defined extensional quotient
object that still bears charged computational information; or (iii) a
covariant compiler whose maps and all costs belong to one realization.
Then prove a cited unconditional lower bound for that same ordinary machine.
Every arrow remains \OPEN and appears separation-strength.

\begin{thebibliography}{9}
\bibitem{Rice53} H. G. Rice, ``Classes of recursively enumerable sets and
their decision problems,'' \emph{Transactions of the American Mathematical
Society} 74(2), 358--366, 1953.
\bibitem{PSS98} A. Pnueli, M. Siegel, and E. Singerman, ``Translation
validation,'' \emph{TACAS 1998}, LNCS 1384, 151--166, 1998.
\bibitem{BGS75} T. Baker, J. Gill, and R. Solovay, ``Relativizations of the
$P=?NP$ question,'' \emph{SIAM Journal on Computing} 4(4), 431--442, 1975.
\bibitem{AW09} S. Aaronson and A. Wigderson, ``Algebrization: a new barrier
in complexity theory,'' \emph{ACM Transactions on Computation Theory} 1(1),
2009.
\bibitem{RR97} A. Razborov and S. Rudich, ``Natural proofs,''
\emph{Journal of Computer and System Sciences} 55(1), 24--35, 1997.
\bibitem{FLMV05} L. Fortnow, R. Lipton, D. van Melkebeek, and A. Viglas,
``Time-space lower bounds for satisfiability,'' \emph{Journal of the ACM}
52(6), 835--865, 2005.
\end{thebibliography}

\noindent\textbf{Final status.}
\claim{TESTED-WRAPPER-MONOID-NONCERTIFICATION1} is \EXACT pending independent
review; the AST audit is \COMP;
\claim{FINITE-WRAPPER-CONTROLS-IMPLY-UNIVERSAL-INVARIANCE1} is \RETIRED.
Universal semantic anti-sharing, arbitrary-decider transfer, a model-matched
endpoint, nonrelativization, nonalgebrization, UCORE-TIME1, and P versus NP
remain \OPEN.

% END_OF_VERIFIED_REPORT::LSPV-3::WRAPORBIT1
\end{document}
